% The data package bobilib/results/: layout, instances.csv, the .solu file,
% solution and certificate formats, checksums. Source: results/README.md.
\section{Data package and formats}\label{app:data}

Every result in the paper is backed by a file in the data package, the
folder \texttt{bobilib/results/}. The package's index,
\texttt{results/README.md}, follows the repository style of the INFORMS
Journal on Computing (cite, description, data, file formats, verification,
replication). The index, \texttt{instances.csv},
\texttt{bobilib\_update.solu} and \texttt{SHA256SUMS} are written by
\texttt{paper/scripts/build\_data\_package.py}; \texttt{make data} in
\texttt{bobilib/paper/} regenerates them, together with the macros of every
number quoted in the paper and the tables of Section~\ref{app:tables}. BOBILib's
statuses are those of its metadata snapshot of 18~October~2026.

\subsection{Layout}

\begin{center}
\small
\begin{tabular}{@{}p{1.5in}p{4.8in}@{}}
\toprule
Path & Contents\\
\midrule
\texttt{instances.csv} & one row per BOBILib instance (\NumInstances{} rows); Section~\ref{app:data-csv}\\
\texttt{bobilib\_update.solu} & status changes and improved values in the MIPLIB/SCIP \texttt{.solu} format; Section~\ref{app:data-solu}\\
\texttt{SHA256SUMS} & SHA-256 checksum of every other file of the package\\
\texttt{summary.json} & per-instance outcome and per-family counts\\
\texttt{kip/} & \NumKipFiles{} knapsack interdiction solution files at the published optima of Weninger and Fukasawa\cite{weninger2025fast}\\
\texttt{ss/} & the \NumNewKip{} subset-sum optima proven here; \texttt{proofs.csv} (optimum, nodes, CPU seconds, nodes of the independent prover)\\
\texttt{or/} & \NumOrPoints{} verified points for the \NumOr{} OR instances; the bounds of each instance; \texttt{compare.csv} (published bounds, BOBILib's value, ours, gaps)\\
\texttt{clique/} & \NumCliqueSet{} partition certificates, solver logs of the part-size proofs, the independent review\\
\texttt{network/} & \NumVerifyNetwork{} attack-protect trees certified by enumeration; \texttt{summary.csv}; the review\\
\texttt{rays/} & \NumRayFiles{} recession directions (the \NumNoFeasible{} open instances and one that BOBILib lists as infeasible); \texttt{ray\_verify.jsonl} (exact checks)\\
\texttt{general/} & \NumGeneralFiles{} points for MIPLIB-derived instances without coupling rows (nine files over 1\,MB in \texttt{large\_points.tar.gz}); \texttt{hpr\_optimal.csv}\\
\texttt{general-coupled/} & \NumCoupledFiles{} points for MIPLIB-derived instances with coupling rows; \texttt{summary.csv}\\
\bottomrule
\end{tabular}
\end{center}

\noindent Each folder has its own \texttt{README.md} with the method, the
columns of its tables, the commands that produced it and, where one was done,
its independent review.

\subsection{The instance table \texttt{instances.csv}}\label{app:data-csv}

The file is UTF-8, comma separated, with one header row and one row per
BOBILib instance. Its columns:

\begin{center}
\small
\begin{tabular}{@{}p{1.5in}p{4.8in}@{}}
\toprule
Column & Meaning\\
\midrule
\texttt{instance} & BOBILib instance name\\
\texttt{collection} & BOBILib sub-collection (last directory of the instance path), e.g.\ \texttt{kp}, \texttt{or}, \texttt{miplib2017}\\
\texttt{family} & \texttt{knapsack}, \texttt{multidim-knapsack}, \texttt{clique}, \texttt{network}, \texttt{generalized}, \texttt{assignment} (interdiction), \texttt{general-pure}, \texttt{general-mixed}\\
\texttt{bobilib\_status}, \texttt{bobilib\_value} & BOBILib's status and best known objective value (empty: none)\\
\texttt{new\_status} & \texttt{optimal} (proven, in published work or here), \texttt{infeasible\_bilevel} (no bilevel feasible point), \texttt{feasible} (first verified point; BOBILib had none), \texttt{unchanged}, \texttt{bobilib\_optimal}, \texttt{bobilib\_infeasible}\\
\texttt{primal\_value} & best upper-level value backed by a verified solution file (of this package or of BOBILib) or by a published optimum\\
\texttt{dual\_bound} & best proven lower bound (the leader minimises), if any\\
\texttt{source} & citation key of the published result (Table~\ref{tab:source-keys}), \texttt{this work}, or \texttt{ThuraufEtAl2026} (BOBILib's own value)\\
\texttt{method} & how the status or value was obtained\\
\texttt{certificate} & path, relative to \texttt{results/}, of the solution or certificate file; \texttt{ARCHIVE.tar.gz:FILE} names a member of an archive\\
\texttt{verified} & \texttt{yes}: the certificate passed its check (\texttt{bobi verify}, or the exact ray check); \texttt{no}: checked in floating point only (the two \texttt{supportcase42} directions); \texttt{na}: no certificate in this package\\
\texttt{notes} & free text, separated by semicolons\\
\bottomrule
\end{tabular}
\end{center}

\noindent Rows per \texttt{new\_status}: \NumBobiOptimal{} \texttt{bobilib\_optimal},
\NumBobiInfeasible{} \texttt{bobilib\_infeasible}, \NumOptimalNow{} \texttt{optimal},
\NumNoFeasible{} \texttt{infeasible\_bilevel}, \NumFirstPoint{} \texttt{feasible}
and \NumUnchanged{} \texttt{unchanged}. The \NumOptimalNow{} optimal rows are the
\NumLitOptimal{} published optima and the \NumNewOptimal{} proven here; together
with the \NumNoFeasible{} rows without a bilevel feasible point they are the
\NumSettled{} settled instances. The keys of the \texttt{source} column map to
the references as follows.

\begin{table}[ht]
\centering
\small
\caption{Citation keys of the \texttt{source} column of \texttt{instances.csv}.}\label{tab:source-keys}
\begin{tabular}{@{}lll@{}}
\toprule
Key in \texttt{source} & Reference & Optima\\
\midrule
\texttt{ThuraufEtAl2026} & BOBILib's own value\cite{thurauf2026bobilib} & \\
\texttt{WeningerFukasawa2025} & Weninger and Fukasawa\cite{weninger2025fast} & \NumLitKip{}\\
\texttt{LozanoBergmanCire2026} & Lozano, Bergman and Cire\cite{lozano2026network} & \NumLitLozano{}\\
\texttt{FischettiEtAl2019} & Fischetti, Ljubi\'c, Monaci and Sinnl\cite{fischetti2019interdiction} & \NumLitImkp{}\\
\texttt{FischettiEtAl2017} & Fischetti, Ljubi\'c, Monaci and Sinnl\cite{fischetti2017new} & \NumLitFischettiSeventeen{}\\
\texttt{VasquezLozanoVanHoeve2025} & V\'asquez, Lozano and van Hoeve\cite{vasquez2025single} & \NumLitVasquez{}\\
\texttt{BOBILibEarlierRelease} & an earlier release of BOBILib's metadata\cite{bobilib_website} & \NumLitEarlierRelease{}\\
\texttt{this work} & a certificate in this package & \NumNewOptimal{}\\
\bottomrule
\end{tabular}
\end{table}

\subsection{The status file \texttt{bobilib\_update.solu}}\label{app:data-solu}

The file follows the MIPLIB/SCIP convention, one line per instance whose
status differs from BOBILib's metadata or whose best value improves:
\begin{center}
\small
\begin{tabular}{@{}lll@{}}
\texttt{=opt=} & \emph{name value} & proven optimal value (published, or proven here)\\
\texttt{=inf=} & \emph{name} & no bilevel feasible point (the follower is unbounded for every leader choice)\\
\texttt{=best=} & \emph{name value} & best verified bilevel feasible point (a solution file here), not proven optimal\\
\texttt{=bestdual=} & \emph{name value} & best proven lower bound, where this work proves a better one than any published\\
\end{tabular}
\end{center}
It has \NumSoluOpt{} \texttt{=opt=} lines (\NumLitOptimal{} published,
\NumNewOptimal{} from this work), \NumSoluInf{} \texttt{=inf=}, \NumSoluBest{}
\texttt{=best=} and \NumSoluBestDual{} \texttt{=bestdual=} lines. The
\texttt{=bestdual=} lines are the \NumOrBetterLb{} OR instances with one
follower knapsack and the \NumCliqueOpen{} open clique instances; an instance
can have both a \texttt{=best=} and a \texttt{=bestdual=} line. Every value is
the upper-level objective $c^{\top}x+d^{\top}y$ of the minimising leader under
the optimistic convention, printed at full precision (integers without a
decimal point, other values as the shortest decimal that reads back as the
same double). Comment lines start with \texttt{\#}.

\subsection{Solution files}

Solution files (\texttt{*.res.json}) are in BOBILib's own format, so BOBILib's
checker and \texttt{bobi verify} read them unchanged. Each is a JSON object
with the fields \texttt{name}, \texttt{bilevel\_type} (always
\texttt{optimistic}), \texttt{status} (\texttt{optimal} or \texttt{feasible}),
\texttt{difficulty}, \texttt{objective\_value} (the upper-level objective,
which BOBILib minimises), and \texttt{upper\_level\_decisions} and
\texttt{lower\_level\_decisions}, which map every variable name of the
instance to its value. A file states \texttt{optimal} only when the instance's
row in \texttt{instances.csv} has \texttt{new\_status} \texttt{optimal}.

\subsection{Certificates}

\paragraph{Recession directions.} \texttt{rays/*.ray.json} maps follower
variable names to the entries of a direction $d$ (its support only), as
decimal numbers; \texttt{*.exact.ray.json} holds exact rationals written as
\texttt{"p/q"}. The certificate is valid if $d$ satisfies the conditions of
Theorem~\ref{thm:ray}: for every follower row, with the leader terms and
right-hand side dropped, $a^{\top}d\le 0$ for a row $a^{\top}y\le\cdot$,
$a^{\top}d\ge 0$ for a row $a^{\top}y\ge\cdot$ and $a^{\top}d=0$ for an
equation; $d_j\ge 0$ for a finite lower bound and $d_j\le 0$ for a finite upper
bound; and $f^{\top}d<0$. \texttt{scripts/verify\_ray.py} checks this in rational
arithmetic and writes \texttt{rays/ray\_verify.jsonl}.

\paragraph{Partitions.} The clique solution files are the certificates: the
leader deletes exactly the edges inside the parts of an $r$-partition (at most
$k$ of them), so the remainder is $r$-partite
(Proposition~\ref{prop:clique-partition}), and the follower part is a maximum
clique of the remainder. \texttt{clique/summary.csv} gives per instance $n$, the
density, $m$, $k$, BOBILib's value, the lower bound with its argument (Mantel,
or Brouwer's theorem with the part-size proofs) and the upper bound with its
partition. The logs of the solver runs behind the \NumCliquePartCases{}
part-size proofs are in \texttt{clique/logs/lb\_solver\_logs.tar.gz}.

\paragraph{Subset-sum proofs.} A proof of $\mathrm{OPT}\ge\ell$ is a
computation, not a file. \texttt{ss/proofs.csv} records, per instance, the
optimum, the nodes and CPU seconds of the search and the nodes of the
independent prover \texttt{scripts/ss\_check.py}, which re-proves the bound
from the instance files (Section~\ref{app:verification}).

\paragraph{Enumerations and HPR bounds.} \texttt{network/summary.csv} gives
per tree the attacks enumerated, the optimum, the number of optimal attacks
and the cross-checks (Theorem~\ref{thm:fire-cert}).
\texttt{general/hpr\_optimal.csv} lists the points that meet the MIPLIB
optimum of the HPR; \texttt{scripts/compare\_hpr.py} re-checks that the HPR is
the MIPLIB model (Lemma~\ref{lem:hpr}).

\paragraph{OR bounds.} \texttt{or/bounds\_one\_knapsack.jsonl} gives, per
instance with one follower knapsack, the lower bound with its two parts (merged-row
online bound, Theorem~\ref{thm:mkip-omega}, and LP bracket,
Theorem~\ref{thm:mkip-bracket}) and the upper bound;
\texttt{or/points\_several\_knapsacks.jsonl} gives, per instance with several
follower knapsacks, the bracket minimiser's LP bound and the certified
follower value.

\subsection{Checksums and how to check the package}

\texttt{SHA256SUMS} has one line per file of the package other than itself,
sorted by path, in the format of \texttt{sha256sum}. With BOBILib's instance
files in \texttt{data/collection} (from \texttt{bobilib.org}) and the
\texttt{bobi} crate built, the following commands, run from \texttt{bobilib/},
repeat the checks of Section~\ref{app:verification}:
\begin{center}
\begin{minipage}{0.96\linewidth}
\footnotesize
\begin{verbatim}
(cd results && sha256sum -c SHA256SUMS)
D=data/collection
./target/release/bobi verify --instances $D results/kip results/ss results/network results/or
./target/release/bobi verify --instances $D --generic results/clique results/general \
    results/general-coupled
python3 -I scripts/verify_ray.py INSTANCE.mps.gz INSTANCE.aux results/rays/INSTANCE.ray.json
python3 -I scripts/ss_check.py INSTANCE.mps.gz INSTANCE.aux - OPTIMUM
python3 -I scripts/compare_hpr.py INSTANCE.mps.gz INSTANCE.aux MIPLIB.mps.gz
\end{verbatim}
\end{minipage}
\end{center}
The nine large general points are checked after unpacking
\texttt{general/large\_points.tar.gz}.
